% Comments

\newcommand{\simon}[1]{%
  {\scriptsize
      \textbf{\textcolor{red}{Simon: #1}}
  }%
}%

\newcommand{\peter}[1]{%
  {\scriptsize
      \textbf{\textcolor{blue}{Peter: #1}}
  }%
}%

% Paper specific notation

\setlength{\intextsep}{10pt}

\crefname{problem}{problem}{problems}

\newcommand{\inputV}{v}
\newcommand{\consistent}{\ensuremath{\mathit{Consistent}}}
\newcommand{\remaining}{\ensuremath{\mathit{Remaining}}}
\newcommand{\dontcare}{\_}
\newcommand{\defined}{\ensuremath{\mathit{defined}}}
\newcommand{\undefined}{\ensuremath{\mathit{undefined}}}
\newcommand{\properties}{P}
\newcommand{\satisfies}{\vDash}
\newcommand{\simulator}{\mathcal{A}}
\newcommand{\Induced}[2]{\llbracket #1 \rrbracket_{#2}}
\newcommand{\timebase}{\setreal_{\geq 0}}
\newcommand{\stepbase}{\setreal_{> 0}}

\newcommand{\stateset}[1]{S_{#1}}
\newcommand{\runstate}[1]{S^{R}_{#1}}
\newcommand{\state}[1]{s_{#1}}
\newcommand{\inputs}[1]{U_{#1}}
\newcommand{\inputvar}[1]{u_{#1}}
\newcommand{\outputs}[1]{Y_{#1}}
\newcommand{\outputvar}[1]{y_{#1}}
\newcommand{\mayReject}{M}
\newcommand{\backtrack}{B}
\newcommand{\values}{\mathcal{V}}
\newcommand{\valuesExchanged}{\mathcal{V_{E}}}
\newcommand{\true}{\mathit{true}}
\newcommand{\false}{\mathit{false}}
\newcommand{\feedthrough}[1]{F_{#1}}
\newcommand{\reactivity}[1]{R_{#1}}
\newcommand{\fset}[1]{\mathtt{set}_{#1}}
\newcommand{\fget}[1]{\mathtt{get}_{#1}}
\newcommand{\fdoStep}[1]{\mathtt{step}_{#1}}
\newcommand{\fpreset}[1]{\mathtt{preSet}_{#1}}
\newcommand{\fpreget}[1]{\mathtt{preGet}_{#1}}
\newcommand{\fpredoStep}[1]{\mathtt{preStep}_{#1}}
\newcommand{\ftime}{\mathtt{ftime}}
\newcommand{\maxStep}[1]{\mathtt{h}_{#1}}
\newcommand{\timestamp}[1]{\varphi(#1)}
\newcommand{\feedsto}[2]{U_{#1}^{#2}}
\newcommand{\CheckCon}{\mathit{CheckConv}}
\newcommand{\stepfound}{\mathit{Step\_found}}
\newcommand{\configuration}{\mathit{Configurations}}


\newcommand{\dom}[1]{\mathit{dom(#1)}}
\newcommand{\algebraic}[1]{\mathit{algebraic_{#1}}}

 

\newcommand{\precondition}{\mathit{Pre}}


\newcommand{\complexity}{Com}
\newcommand{\master}{\mathcal{A}}
\newcommand{\alloutputs}{Y}
\newcommand{\allfeedthroughs}{F}
\newcommand{\allreactivity}{R}
\newcommand{\alldelayed}{D}

\newcommand{\allcontracts}{\mathcal{C}}
\newcommand{\coupling}{L}
\newcommand{\allinputs}{U}
\newcommand{\fmus}{C}
\newcommand{\fmusLoop}{C_{loop}}
\newcommand{\sequence}[1]{\pargroup{#1}}
\newcommand{\functioncall}{f}
\newcommand{\initcall}{I}
\newcommand{\allfunctioncalls}{F}
\newcommand{\fmu}[1]{\texttt{#1}}
\newcommand{\signal}[1]{\texttt{#1}}
\newcommand{\before}[2]{\ensuremath{#1 \twoheadrightarrow #2}}
\newcommand{\ibefore}[2]{\ensuremath{#1 \rightarrow #2}}
\newcommand{\after}[1]{{#1}'}
\newcommand{\aftern}[2]{{#1}^{(#2)}}
\newcommand{\stateafter}[2]{\ensuremath{\state{#1}^{(#2)}}}
\newcommand{\minstep}{\ensuremath{h_{min}}}


\newcommand{\etime}{\ensuremath{E}}

\newcommand{\verifier}{\mathit{Verifier}}

\newcommand{\turnOn}{\mathit{TurnPreconditionsOn()}}
\newcommand{\turnOff}{\mathit{TurnPreconditionsOff()}}

\newtheorem{procedure}{Procedure}{}
\newtheorem{assumption}{Assumption}{}

\newtheorem{experiment}{Experiment}{}

\algrenewcommand\alglinenumber[1]{\scriptsize #1:}% Adapt size of line numbers in algorithm

\algnewcommand\algorithmicforeach{\textbf{for each}}
\algdef{S}[FOR]{ForEach}[1]{\algorithmicforeach\ #1\ \algorithmicdo}
% New definitions
\algnewcommand\algorithmicswitch{\textbf{switch}}
\algnewcommand\algorithmiccase{\textbf{case}}
\algnewcommand\algorithmicassert{\texttt{assert}}
\algnewcommand\Assert[1]{\State \algorithmicassert(#1)}%
% New "environments"
\algdef{SE}[SWITCH]{Switch}{EndSwitch}[1]{\algorithmicswitch\ #1\ \algorithmicdo}{\algorithmicend\ \algorithmicswitch}%
\algdef{SE}[CASE]{Case}{EndCase}[1]{\algorithmiccase\ #1}{\algorithmicend\ \algorithmiccase}%
\algtext*{EndSwitch}%
\algtext*{EndCase}%



% Generic stuff

\newcommand{\footurl}[1]{\footnote{\url{#1}}}

% Math
\newcommand{\brackets}[1]{\ensuremath{ \left[ #1 \right] }}
\newcommand{\tuple}[1]{\ensuremath{ \left\langle #1 \right\rangle }}
\newcommand{\set}[1]{\ensuremath{ \left\{ #1 \right\}}}
\newcommand{\system}[1]{\ensuremath{ \begin{cases} #1 \end{cases}}}
\newcommand{\rightgroup}[1]{\ensuremath{ \left. \begin{matrix} #1 \end{matrix} \right\} } }
\newcommand{\pargroup}[1]{\ensuremath{ \left( #1 \right)}}
\newcommand{\inv}[1]{\ensuremath{\pargroup{ #1 }^{-1}}}
\newcommand{\dert}[1]{\ensuremath{ \dot{#1} }}
\newcommand{\ddert}[1]{\ensuremath{ \ddot{#1} }}
\newcommand{\partialder}[2]{\ensuremath{ \frac{\partial#1}{\partial#2} }}
\newcommand{\setreal}[0]{\ensuremath{\mathbb{R}}}
%\newcommand{\setbool}[0]{\ensuremath{\mathit{Bool}}}
\newcommand{\setnat}[0]{\ensuremath{\mathbb{N}}}
\newcommand{\norm}[1]{\left\lVert#1\right\rVert}
\newcommand{\bnorm}[1]{\big\lVert#1\big\rVert}
\newcommand{\abs}[1]{\left|#1\right\|}
\newcommand{\xs}[2]{\ensuremath{#1^{\left[#2\right]}}}
\newcommand{\infinitynorm}[1]{\left\lVert#1\right\rVert_\infty}


\newcommand{\vectorOne}[1]{\brackets{%
\begin{matrix}
  #1
 \end{matrix}%
}}
\newcommand{\vectorTwo}[2]{\brackets{%
\begin{matrix}
  #1 \\
  #2
 \end{matrix}%
}}
\newcommand{\vectorThree}[3]{\brackets{%
\begin{matrix}
  #1 \\
  #2 \\
  #3
 \end{matrix}%
}}
\newcommand{\vectorFour}[4]{\brackets{%
\begin{matrix}
  #1 \\
  #2 \\
  #3 \\
  #4
 \end{matrix}%
}}
\newcommand{\vectorFive}[5]{\brackets{%
\begin{matrix}
  #1 \\
  #2 \\
  #3 \\
  #4 \\
  #5
 \end{matrix}%
}}
\newcommand{\vectorSix}[6]{\brackets{%
\begin{matrix}
  #1 \\
  #2 \\
  #3 \\
  #4 \\
  #5 \\
  #6
 \end{matrix}%
}}
\newcommand{\vectorSeven}[7]{\brackets{%
\begin{matrix}
  #1 \\
  #2 \\
  #3 \\
  #4 \\
  #5 \\
  #6 \\
  #7
 \end{matrix}%
}}
\newcommand{\vectorEight}[8]{\brackets{%
\begin{matrix}
  #1 \\
  #2 \\
  #3 \\
  #4 \\
  #5 \\
  #6 \\
  #7 \\
  #8
 \end{matrix}%
}}

\newenvironment{aligneq*}%
{
\begin{equation*}
\begin{aligned}
}{
\end{aligned}
\end{equation*}
}

\newenvironment{aligneq}%
{
\begin{equation}
\begin{aligned}
}{
\end{aligned}
\end{equation}
}



%enable \cref{...} and \Cref{...} instead of \ref: Type of reference included in the link
%Nice formats for \cref
\usepackage{iflang}
\IfLanguageName{ngerman}{
  \crefname{table}{Tab.}{Tab.}
  \Crefname{table}{Tabelle}{Tabellen}
  \crefname{figure}{\figurename}{\figurename}
  \Crefname{figure}{Abbildungen}{Abbildungen}
  \crefname{equation}{Gleichung}{Gleichungen}
  \Crefname{equation}{Gleichung}{Gleichungen}
  \crefname{listing}{\lstlistingname}{\lstlistingname}
  \Crefname{listing}{Listing}{Listings}
  \crefname{section}{Abschnitt}{Abschnitte}
  \Crefname{section}{Abschnitt}{Abschnitte}
  \crefname{paragraph}{Abschnitt}{Abschnitte}
  \Crefname{paragraph}{Abschnitt}{Abschnitte}
  \crefname{subparagraph}{Abschnitt}{Abschnitte}
  \Crefname{subparagraph}{Abschnitt}{Abschnitte}
}{
  \crefname{section}{Sect.}{Sect.}
  \Crefname{section}{Section}{Sections}
  \crefname{listing}{\lstlistingname}{\lstlistingname}
  \Crefname{listing}{Listing}{Listings}
}
\crefname{section}{Sect.}{Sect.}
\Crefname{section}{Section}{Sections}
\crefname{listing}{\lstlistingname}{\lstlistingname}
\Crefname{listing}{Listing}{Listings}

\crefname{assumption}{Assumption}{Assumptions}


% update later with official instructions from Springer
\def\orcidID#1{${}^{\smash{\href{https://orcid.org/#1}{\protect\raisebox{-1.25pt}{\protect\includegraphics{images/orcid_color-eps-converted-to}}}}}$}
